\documentclass[11pt,a4paper]{article}
\usepackage[margin=1in]{geometry}
\usepackage[T1]{fontenc}
\usepackage[utf8]{inputenc}
\usepackage{lmodern}
\usepackage{microtype}
\usepackage{enumitem}
\usepackage{xcolor}
\usepackage{hyperref}
\hypersetup{colorlinks=true, linkcolor=black, urlcolor=blue}
\setlength{\parindent}{0pt}
\setlength{\parskip}{0.55em}
\setlist[itemize]{leftmargin=1.5em,itemsep=0.2em,topsep=0.2em}
\setlist[enumerate]{leftmargin=1.7em,itemsep=0.2em,topsep=0.2em}
\pagenumbering{arabic}
\begin{document}
\begin{center}
{\LARGE\bfseries Swimming Performance Platform Proposal}
\end{center}
\vspace{0.5em}

Hello,\par

My name is Fabricio Santana. I am a software engineer and technical lead with 20+ years of experience in backend development, software architecture, database integration, automated testing, cloud computing, and delivery of critical systems. I have also led large engineering teams and worked on projects where data consistency, business rules, validation, security, and maintainability are essential.\par

Your project is a strong match for my background because the main challenge is not only importing Excel files. The important part is building a reliable data import workflow that protects the integrity of swimmer profiles, competition results, scores, and rankings.\par


\section*{Working methodology}

My delivery method is designed to reduce risk before development starts. The work begins with a fixed Discovery Sprint and then moves into a delivery package only after scope, priorities, responsibilities, dependencies, and acceptance criteria are clear.\par

The first step is a one-week Discovery Sprint with a fixed price of USD 500. In this phase, I can hold up to five remote meetings to understand the business problem, review the current platform, analyze the Excel files, map requirements, identify risks, review integrations, and define priorities. The deliverable is a practical work plan for the implementation, including recommended scope, technical approach, milestones, success criteria, and the most appropriate delivery package.\par

After discovery, the project can move into one of three fixed-price, four-week delivery packages. The package is selected according to scope, complexity, delivery speed, integration needs, quality requirements, and the amount of engineering support required. Some roles may participate part-time depending on the agreed scope, especially Scrum Master, technical leadership, test automation, DevOps, observability, and deployment support.\par

\textbf{Delivery package options}\par

\begin{itemize}
\item \textbf{Essential Delivery} --- USD 2,000, 4 weeks. Includes a Scrum Master and one full-stack engineer. This package is designed for a focused module, MVP, backend feature, import flow, or small integration with limited uncertainty. It is suitable when the scope is well defined, the number of integrations is limited, and the main goal is to deliver a reliable working feature without unnecessary overhead.
\item \textbf{Product Acceleration} --- USD 4,000, 4 weeks. Includes a Scrum Master, two full-stack engineers, and one engineer focused on test automation, DevOps, observability, and deployment. This package is recommended when faster delivery is needed or when the scope includes a product, API, integration, internal workflow, multiple features, or more demanding quality needs.
\item \textbf{Scale \& Intelligence} --- USD 6,000, 4 weeks. Includes a Technical lead, Scrum Master, two full-stack engineers, and one engineer focused on test automation, DevOps, observability, and deployment. This package is designed for complex systems, advanced data workflows, automation, auditability, security, observability, broader platform evolution, or stronger technical scale planning.
\end{itemize}

Role responsibilities are clear: the Scrum Master organizes scope, communication, ceremonies, progress tracking, and delivery alignment; full-stack engineers implement backend, frontend, API, and database changes as needed; the test automation, DevOps, observability, and deployment engineer supports automated tests, CI/CD, monitoring, release quality, and deployment; and the Technical lead, included in the Scale \& Intelligence package, owns architecture, technical decisions, code quality, and risk management for more complex projects.\par

Each delivery cycle is executed with Scrum-inspired practices, frequent alignment, technical review, automated testing when applicable, and a demonstration of what was delivered. The client keeps ownership of the produced code, and the scope is agreed before execution to avoid open-ended work and unclear expectations.\par

\textbf{Construction defect warranty.} After delivery and acceptance, I include a 30-day warranty for construction defects related to the approved scope. This covers fixes for implementation errors, broken agreed behavior, or defects introduced by the delivered code. It does not cover new features, changes in business rules, new Excel formats not included in the approved scope, third-party service changes, infrastructure changes outside the project, or requests that expand the original acceptance criteria.\par

\section*{Proposed work plan}

I would approach the project in structured phases:\par

\begin{enumerate}
\item Technical assessment and import design
\end{enumerate}

\begin{itemize}
\item Review the current application architecture, database schema, existing import logic, and ranking calculation rules.
\item Analyze representative Excel files with different layouts.
\item Define the import workflow, validation rules, duplicate detection strategy, and matching criteria.
\item Confirm the acceptance criteria before implementation.
\end{itemize}

\begin{enumerate}
\item Excel import and validation process
\end{enumerate}

\begin{itemize}
\item Implement or improve Excel/CSV parsing for the expected file structures.
\item Normalize imported data into a controlled staging process before saving final records.
\item Validate required fields, formats, competition metadata, events, swimmer data, times/results, and category information.
\item Generate clear validation feedback for rows that cannot be safely imported.
\item Prevent duplicate entries using deterministic rules based on competition, swimmer, event, date, result, and source data.
\end{itemize}

\begin{enumerate}
\item Swimmer matching and reconciliation
\end{enumerate}

\begin{itemize}
\item Prioritize exact matching using swimmer IDs or other unique identifiers when available.
\item Add name-based fuzzy matching only as a fallback.
\item Use confidence levels to avoid unsafe automatic matches.
\item Flag unmatched or ambiguous swimmers for manual review before final reconciliation.
\end{itemize}

\begin{enumerate}
\item Score and ranking calculation updates
\end{enumerate}

\begin{itemize}
\item Review the existing point and ranking formulas.
\item Implement the updated scoring criteria in a clear and testable way.
\item Define when scores and rankings should be recalculated after each import.
\item Recalculate only the affected rankings when possible, while preserving correctness.
\item Add automated tests for scoring and ranking edge cases.
\end{itemize}

\begin{enumerate}
\item Testing, documentation, and handoff
\end{enumerate}

\begin{itemize}
\item Test the import process with multiple Excel examples and edge cases.
\item Add technical documentation for the import flow, matching logic, duplicate prevention, and ranking recalculation.
\item Provide a handoff session or written notes so your team can maintain and extend the feature.
\end{itemize}

\section*{Estimated timeline}

For the recommended starting scope, my estimated timeline is 5 weeks after receiving repository access, database/schema information, sample Excel files, and the current scoring rules: 1 week for Discovery Sprint and 4 weeks for one Essential Delivery cycle.\par

A practical schedule would be:\par

\begin{itemize}
\item Week 1: Discovery Sprint, technical assessment, import design, and confirmation of rules.
\item Weeks 2 and 3: Excel import, validation, duplicate prevention, and swimmer matching.
\item Weeks 4 and 5: scoring/ranking updates, recalculation logic, and automated tests.
\item Optional extension, only if needed: refinements, additional edge cases, broader documentation, or scope items identified during discovery.
\end{itemize}

If discovery confirms that the current assumptions are correct, one Essential Delivery cycle should be enough. If the existing codebase, Excel formats, or ranking rules are more complex than expected, the timeline and package recommendation should be adjusted before implementation starts.\par

\section*{Availability and communication}

I can work remotely and keep communication in English through Workana messages, scheduled calls, and written progress updates. I usually prefer short alignment checkpoints during the week, especially after reviewing the codebase and after each major implementation milestone.\par

This helps reduce misunderstandings and keeps the project aligned with your actual data, business rules, and operational workflow.\par

\section*{Budget}

Based on the current understanding, my recommended starting budget is USD 2,500: USD 500 for the one-week Discovery Sprint plus USD 2,000 for one Essential Delivery cycle, assuming:\par

\begin{itemize}
\item The platform already exists and the work is focused on improving the import, matching, scoring, and ranking logic.
\item You can provide access to the codebase, database schema, and anonymized Excel samples.
\item The number of Excel formats is limited and can be mapped during the first phase.
\item The updated scoring criteria can be clearly defined with your input.
\end{itemize}

The Discovery Sprint will refine the scope and confirm whether one Essential Delivery cycle is sufficient. If the review shows many unknown Excel formats, major database refactoring, a new admin interface, complex historical data cleanup, or advanced scoring changes, I will recommend moving to Product Acceleration or Scale \& Intelligence before implementation starts.\par

\section*{Questions before final confirmation}

Before confirming the final scope, I would like to clarify a few points:\par

\begin{enumerate}
\item What technology stack is the current platform using for the backend and database?
\item Can you provide anonymized examples of the Excel files, including the different structures you receive?
\item Do most competition files include a unique swimmer ID, or is name-based matching frequently required?
\item What fields should define a duplicate result: swimmer, event, competition, date, time/result, category, or another combination?
\item Do you already have documented rules for the points and ranking calculations?
\item Should valid rows be imported automatically while unmatched rows are held for review, or should the entire import wait for manual approval?
\item Is there already an admin screen for reviewing unmatched swimmers, or should this be added as part of the project?
\item Should ranking recalculation apply to all historical results or only affected competitions, seasons, age groups, or categories?
\item Does the project currently have automated tests or a staging environment?
\item Are there any deadlines related to upcoming competitions or ranking publications?
\end{enumerate}

Once these points are clear, I can confirm the final delivery plan and milestones.\par

\end{document}
